\Article{11}
Imagine straight lines drawn from each integer point $k$ to $k+\frac12i$, their right sides regarded as positive and their left sides as negative. Let $l$ be a line, and let $Sl$ denote the sum of all intersections of $l$ with this system of lines, counting positively those where $l$ passes from the negative side to the positive, and negatively the opposite intersections. The decident then has the following expression:
\[
\Sigma(\Upsilon l\pdot Sl)+\Sigma Sl'-S\mu
\]

Here $l$ denotes all ligatures of the squares in $\omega$ (always taken so that the squares lie to their right), and $l'$ denotes the ligatures on the boundary of the figure $\omega$ between $0$ and $\frac12m$, hence outside $\Omega'$.

All ligatures $l$ consist of
\begin{enumerate}
\item[1)] $l'$;
\item[2)] $l''$, the boundary ligatures inside $\Omega'$, hence $\lambda',\overline{\overline\nu},\overline{\overline\lambda}$;
\item[3)] $l'''$, which lie in the interior of $\omega$.
\end{enumerate}
\Needspace{27\baselineskip}
If $l$ were allowed to range over all ligatures inside $\omega$ or on the boundary of this figure, insofar as they start at points $\frac{km}{2M}$ such that $k$ is divisible by $1+i$, the decident would be
\[
=\Sigma\alpha\pdot Sl-S\mu
\]

\begin{list}{}{\setlength{\labelwidth}{5em}\setlength{\leftmargin}{5.5em}}
\item[where $\alpha=1$] for all ligatures in the interior of $\omega$;
\item[$\alpha=\Upsilon l+1$] for all boundary ligatures outside $\Omega'$ whose direction follows the boundary running from $0$ to $\frac12m$;
\end{list}

\SourceLeaf{341}

\begin{list}{}{\setlength{\labelwidth}{9em}\setlength{\leftmargin}{9.5em}}
\item[$\alpha=-(\Upsilon\overline{\overline\lambda}+1)=-\Upsilon l$] for all those on this boundary running in the opposite direction;
\item[$\alpha=\Upsilon l$] for all boundary ligatures inside $\Omega'$ directed towards $0$;
\item[$\alpha=-\Upsilon l+1$] for all boundary ligatures inside $\Omega'$ directed away from $0$.
\end{list}

\Article{12}
We can now combine all the $l$ that occur (in the last manner) in pairs, namely $l$ with $\overline{\overline{\frac12m-l}}$, which we shall call associated ligatures. A single one is excepted, standing isolated or identical with its associated ligature, namely the one running from
\[
\tfrac12(M-1)\pdot\frac{m}{2M}\quad\text{to}\quad\tfrac12(M+1)\frac{m}{2M}
\]
For associated ligatures, $\alpha$ is always the same.
\BookRule{.3\textwidth}
